\documentclass[11pt,a4paper]{article}

\usepackage[margin=1in]{geometry}
\usepackage[utf8]{inputenc}
\usepackage[T1]{fontenc}
\usepackage{hyperref}
\usepackage{amsmath}
\usepackage{amssymb}
\usepackage{amsfonts}
\usepackage{booktabs}
\usepackage{xcolor}
\usepackage{microtype}
\usepackage{enumitem}
\usepackage{amsthm}
\usepackage{mdframed}

% ---- macros ----
\newcommand{\E}{\mathbb{E}}
\newcommand{\Var}{\mathrm{Var}}
\newcommand{\neff}{n_{\mathrm{eff}}}
\newcommand{\ghat}{\hat{g}_t}
\newcommand{\gtrue}{g_t}
\newcommand{\pith}{\pi_{\theta_t}}
\newcommand{\alphaopt}{\alpha^*(t)}

\newcommand{\gap}[1]{\textcolor{red}{\textbf{[GAP: #1]}}}
\newcommand{\claim}[1]{\textcolor{blue}{\textit{#1}}}

% ---- theorem environments ----
\newtheorem{definition}{Definition}
\newtheorem{proposition}{Proposition}
\newtheorem{remark}{Remark}
\newtheorem{assumption}{Assumption}
\newtheorem{lemma}{Lemma}

% framed box for the research question
\newmdenv[linecolor=black,linewidth=1.5pt,innertopmargin=8pt,innerbottommargin=8pt]{rqbox}

% =========================================================
\title{Theoretical Framework:\\
       Effective Sample Degradation in On-Policy RL\\
       and Its Implications for Update Rule Design}

\date{2026-04-19 (draft)}

\begin{document}
\maketitle

\begin{abstract}
We derive, without logical gaps, the following chain:
\emph{on-policy RL training $\to$ policy concentrates over time $\to$ effective sample size
$\neff(t)$ decreases $\to$ optimal learning rate $\alphaopt$ decreases linearly with $\neff(t)$
$\to$ fixed LR eventually exceeds $\alphaopt$ $\to$ gradient noise dominates signal
$\to$ unreliable updates degrade training path quality $\to$ suboptimal final performance.}
This chain identifies a precise, well-defined research gap and motivates a family of solutions
that go beyond learning-rate scheduling.
Each step is derived explicitly; no step is ``by standard result'' without proof.
\end{abstract}

\tableofcontents
\newpage

% =========================================================
\section{Setup}

\paragraph{Objective.}
The optimization target is the expected return under the current policy:
\begin{equation}
  J(\theta) = \E_{\tau \sim \pi_\theta}[R(\tau)]
  \label{eq:objective}
\end{equation}
The reward function $R$, the environment, and the optimal policy $\pi^*$ are all \emph{fixed}.
This is a stationary optimization problem in policy space.

\paragraph{True gradient.}
By the policy gradient theorem:
\begin{equation}
  \gtrue = \nabla_\theta J(\theta_t)
           = \E_{\tau \sim \pith}\!\left[A_t(\tau)\,\nabla_\theta \log \pith(\tau)\right]
  \label{eq:true-gradient}
\end{equation}
where $A_t(\tau) = R(\tau) - b_t$ is the advantage with baseline $b_t$.
This expectation is over \emph{all} possible trajectories, each weighted by $\pith(\tau)$.

\paragraph{Practical estimator.}
$\gtrue$ cannot be computed exactly.
With $n$ rollouts sampled from $\pith$, we use the Monte Carlo estimator:
\begin{equation}
  \ghat = \frac{1}{n}\sum_{i=1}^{n} A_i\,\nabla_\theta \log \pith(\tau_i),
  \qquad \tau_i \overset{\mathrm{i.i.d.}}{\sim} \pith
  \label{eq:estimator}
\end{equation}
Because $\E[\ghat] = \E[A\nabla\log\pith] = \gtrue$, the estimator is \emph{unbiased}.

\paragraph{The actual parameter update.}
\begin{equation}
  \theta_{t+1} = \theta_t + \alpha\,\ghat
               = \theta_t + \alpha\,\gtrue + \underbrace{\alpha(\ghat - \gtrue)}_{\text{estimation error}}
  \label{eq:update}
\end{equation}
The estimation error $\alpha(\ghat - \gtrue)$ is a random variable with mean zero.
Even though the update is unbiased in expectation, its \emph{variance} matters because
the objective is nonlinear: a large random perturbation can move $\theta$ to a region of
lower $J$ even if its expected direction is correct.

% =========================================================
\section{Gradient Estimation Quality}

\subsection{Mean squared error of the gradient estimator: full derivation}
\label{sec:mse}

Let $X_i := A_i\,\nabla_\theta \log \pith(\tau_i) \in \mathbb{R}^d$ denote the $i$-th sample.
The $X_i$ are i.i.d.\ copies of the random vector $X := A\,\nabla\log\pith(\tau)$ with $\tau\sim\pith$.

\textbf{Step 1: Express ĝ − g as a sample mean of centered variables.}

Since $\E[X_i] = \gtrue$, we have $X_i - \gtrue$ is mean-zero, so:
\begin{equation}
  \ghat - \gtrue
  = \frac{1}{n}\sum_{i=1}^{n} X_i - \gtrue
  = \frac{1}{n}\sum_{i=1}^{n} (X_i - \gtrue)
  \label{eq:centered}
\end{equation}

\textbf{Step 2: Expand the squared norm.}

\begin{align}
  \E\!\left[\|\ghat - \gtrue\|^2\right]
  &= \E\!\left[\left\|\frac{1}{n}\sum_{i=1}^{n}(X_i - \gtrue)\right\|^2\right] \notag\\
  &= \frac{1}{n^2}\,\E\!\left[\left\|\sum_{i=1}^{n}(X_i - \gtrue)\right\|^2\right] \notag\\
  &= \frac{1}{n^2}\,\E\!\left[\sum_{i=1}^{n}\sum_{j=1}^{n}
       (X_i - \gtrue)^\top(X_j - \gtrue)\right]
  \label{eq:expand}
\end{align}

\textbf{Step 3: Use i.i.d.\ to eliminate cross terms.}

For $i \neq j$, the variables $X_i - \gtrue$ and $X_j - \gtrue$ are independent and each has mean zero.
Therefore:
\begin{equation}
  \E\!\left[(X_i - \gtrue)^\top(X_j - \gtrue)\right]
  = \E\!\left[X_i - \gtrue\right]^\top \E\!\left[X_j - \gtrue\right]
  = 0^\top \cdot 0 = 0
  \qquad (i \neq j)
  \label{eq:cross-zero}
\end{equation}

Only the diagonal terms $i = j$ survive:
\begin{equation}
  \E\!\left[\|\ghat - \gtrue\|^2\right]
  = \frac{1}{n^2}\sum_{i=1}^{n}\E\!\left[\|X_i - \gtrue\|^2\right]
  = \frac{1}{n^2}\cdot n\cdot\E\!\left[\|X - \gtrue\|^2\right]
  = \frac{1}{n}\,\E\!\left[\|X - \gtrue\|^2\right]
  \label{eq:diagonal}
\end{equation}
where in the last equality we used the fact that all $n$ terms are equal by the i.i.d.\ assumption.

\textbf{Step 4: Identify the remaining expectation as the trace of the covariance matrix.}

Let $\Sigma_t := \mathrm{Cov}_{\pith}[X] = \E\!\left[(X - \gtrue)(X-\gtrue)^\top\right] \in \mathbb{R}^{d\times d}$.
Then:
\begin{equation}
  \E\!\left[\|X - \gtrue\|^2\right]
  = \E\!\left[\mathrm{tr}\!\left((X-\gtrue)(X-\gtrue)^\top\right)\right]
  = \mathrm{tr}\!\left(\E\!\left[(X-\gtrue)(X-\gtrue)^\top\right]\right)
  = \mathrm{tr}(\Sigma_t)
  \label{eq:trace}
\end{equation}
where the second equality uses linearity of trace and linearity of expectation.

\textbf{Conclusion.}

Substituting \eqref{eq:trace} into \eqref{eq:diagonal}:
\begin{equation}
  \boxed{
  \E\!\left[\|\ghat - \gtrue\|^2\right]
  = \frac{\mathrm{tr}(\Sigma_t)}{n}
  }
  \label{eq:mse}
\end{equation}
We define $\sigma_t^2 := \mathrm{tr}(\Sigma_t)$, the total per-sample variance (sum of per-coordinate variances).
Thus the mean squared error of the gradient estimator is $\sigma_t^2 / n$.

\begin{remark}
  Equation~\eqref{eq:mse} holds for any distribution $\pith$, any policy $\theta_t$,
  and any $n \geq 1$.
  It does not yet involve any assumption about policy concentration.
  The factor $n$ in the denominator is exact (not approximate) when all $n$ draws are statistically independent,
  which holds when each rollout is generated by independently sampling from $\pith$.
\end{remark}

\subsection{Effective sample size under policy concentration}
\label{sec:neff}

Equation~\eqref{eq:mse} shows the MSE decreasing as $1/n$.
However, this analysis assumed all $n$ samples are statistically independent.
In practice, when $\pith$ is concentrated (low entropy), many of the $n$ rollouts will be
\emph{identical}, providing no additional information beyond what a single rollout already gives.
We now make this intuition precise.

\subsubsection{Why duplicate samples don't reduce MSE}

Suppose $\pith$ assigns all probability to a single trajectory $\tau^*$: $\pith(\tau^*) = 1$.
Then every rollout produces $\tau_i = \tau^*$, and:
\begin{equation}
  \ghat = \frac{1}{n}\sum_{i=1}^n A^*\nabla\log\pith(\tau^*)
        = A^*\nabla\log\pith(\tau^*)
\end{equation}
This is deterministic (zero variance) but is only one point estimate of $g_t$.
Adding more samples does not change $\ghat$ at all.
From equation~\eqref{eq:mse}, $\mathrm{tr}(\Sigma_t) = 0$ in this case (no variance), so the formula is consistent—but the information content is the same as $n=1$.

Now suppose $\pith$ is concentrated but not fully: it puts probability $1-\epsilon$ on $\tau^*$ and spreads the remaining $\epsilon$ over other trajectories.
With $n=8$ samples, the probability that \emph{all} 8 land on $\tau^*$ is $(1-\epsilon)^8$.
For small $\epsilon$ (e.g., $\epsilon = 0.05$), this probability is $(0.95)^8 \approx 0.66$.
So $66\%$ of the time, all 8 samples are identical, and the effective information is still just 1 sample.

\subsubsection{Distinct trajectories as a measure of information content}

Let $K_n$ denote the number of \emph{distinct} trajectories among $n$ i.i.d.\ draws from $\pith$.

\begin{proposition}
  \label{prop:distinct}
  The expected number of distinct trajectories satisfies:
  \begin{equation}
    \E[K_n] = \sum_{\tau} \left[1 - (1 - \pith(\tau))^n\right]
    \label{eq:expected-K}
  \end{equation}
\end{proposition}

\begin{proof}
  For each trajectory $\tau$, define the indicator $\mathbf{1}[\tau \in \{\tau_1,\ldots,\tau_n\}]$,
  which equals 1 if $\tau$ appears at least once among the $n$ draws, and 0 otherwise.
  By linearity of expectation:
  \begin{equation}
    \E[K_n] = \sum_\tau \E\!\left[\mathbf{1}[\tau \in \{\tau_1,\ldots,\tau_n\}]\right]
             = \sum_\tau \Pr(\tau \in \{\tau_1,\ldots,\tau_n\})
  \end{equation}
  The probability that $\tau$ is \emph{absent} from all $n$ draws is
  $(1 - \pith(\tau))^n$ (by independence and identical distribution),
  so the probability it appears at least once is $1 - (1 - \pith(\tau))^n$, giving \eqref{eq:expected-K}.
\end{proof}

\subsubsection{The approximation $\neff \approx \min(n, \exp(H))$}
\label{sec:neff-approx}

We now derive the approximation.
The entropy of $\pith$ is $H(\pith) = -\sum_\tau \pith(\tau)\log\pith(\tau)$,
and the \emph{perplexity} is $\exp(H(\pith))$.

\textbf{Exact computation for a uniform distribution.}
Suppose $\pith$ is uniform over $M$ trajectories: $\pith(\tau) = 1/M$ for each $\tau$ in the support.
Then $H = \log M$, so $\exp(H) = M$, and by \eqref{eq:expected-K}:
\begin{equation}
  \E[K_n] = M\cdot\left[1 - \left(1 - \frac{1}{M}\right)^n\right]
  \label{eq:uniform-Kn}
\end{equation}
Analyze the two regimes:
\begin{itemize}
  \item \textbf{When $n \ll M$}: By the approximation $(1 - 1/M)^n \approx e^{-n/M} \approx 1 - n/M$
    (valid for $n/M \ll 1$), we get $\E[K_n] \approx M \cdot (n/M) = n$.
    Each new draw lands on a fresh trajectory with high probability.
  \item \textbf{When $n \gg M$}: We have $(1-1/M)^n \approx e^{-n/M} \to 0$,
    so $\E[K_n] \approx M = \exp(H)$.
    Essentially all $M$ trajectories have been seen.
\end{itemize}
Therefore $\E[K_n] \approx \min(n, M) = \min(n, \exp(H))$ for a uniform distribution.

\textbf{Extension to non-uniform distributions.}
For a general distribution, the perplexity $\exp(H)$ represents the effective number of
equally-likely outcomes; it is the size of the uniform distribution with the same entropy.
The approximation $\E[K_n] \approx \min(n, \exp(H))$ holds exactly at the two extremes:
when $\pith$ is perfectly concentrated ($H = 0$, $\exp(H) = 1$, $\E[K_n] = 1$)
and when $\pith$ is uniform ($H = \log M$, $\exp(H) = M$, $\E[K_n] = \min(n, M)$).
For intermediate distributions it is an approximation,
and it overestimates slightly when the distribution has heavy tails
(many low-probability trajectories).

\begin{definition}[Effective sample size]
  \label{def:neff}
  \begin{equation}
    \neff(t) := \E[K_n] \approx \min\!\left(n,\; \exp(H(\pith))\right)
    \label{eq:neff-def}
  \end{equation}
\end{definition}

\textbf{Why this is the right measure of information content.}
Equation~\eqref{eq:mse} assumes $n$ independent samples.
When the policy is concentrated and many samples are duplicates, the effective number of
independent pieces of information is $\neff$, not $n$.
Replacing $n$ with $\neff$ in \eqref{eq:mse} gives the effective MSE:
\begin{equation}
  \E\!\left[\|\ghat - \gtrue\|^2\right]_{\mathrm{eff}}
  \;\approx\;
  \frac{\sigma_t^2}{\neff(t)}
  \label{eq:mse-neff}
\end{equation}

\paragraph{Empirical anchor.}
In our experiments, $n = 8$ rollouts per prompt (GRPO).
At the start of training: $H(\pith) \approx 0.68$ nats, so $\exp(H) \approx 1.97$,
giving $\neff \approx \min(8, 1.97) \approx 2$.
After training: $H(\pith) \approx 0.07$ nats, so $\exp(H) \approx 1.07$,
giving $\neff \approx 1$.
The gradient estimate degrades from roughly 2 effective samples to essentially 1
over the course of training---halving the information in the gradient estimate.

\begin{remark}
  Note: $H = 0.68$ nats is the measured entropy for a 1.5B model at the \emph{response level}.
  The per-token entropy (as logged in experiments) is much higher and is not the quantity used here.
  The response-level entropy measures the concentration of the distribution over complete output sequences.
\end{remark}

% =========================================================
\section{Optimal Learning Rate as a Function of $\neff$}

\subsection{Per-step expected improvement bound: full derivation}
\label{sec:bound-derivation}

We want to find the $\alpha$ that maximizes the expected improvement in $J$ per step.

\textbf{Step 1: Apply the L-smoothness assumption.}

Assume $J(\theta)$ is $L$-smooth, meaning $\nabla J$ is $L$-Lipschitz:
\begin{equation}
  \|\nabla J(\theta') - \nabla J(\theta)\| \leq L\,\|\theta' - \theta\|
  \quad \text{for all } \theta, \theta'
  \label{eq:lipschitz}
\end{equation}
A standard result in optimization (see~\cite{bottou2018optimization}, Lemma A.1) is that $L$-smoothness
implies the following \emph{descent lemma}:
\begin{equation}
  J(\theta') \geq J(\theta) + \nabla J(\theta)^\top(\theta' - \theta) - \frac{L}{2}\|\theta' - \theta\|^2
  \label{eq:descent-lemma}
\end{equation}
for all $\theta, \theta'$.
(Note: for maximization we flip the inequality relative to the standard descent lemma.)

\textbf{Step 2: Substitute the update rule.}

Set $\theta' = \theta_{t+1} = \theta_t + \alpha\ghat$ in \eqref{eq:descent-lemma}:
\begin{equation}
  J(\theta_{t+1})
  \geq J(\theta_t) + \nabla J(\theta_t)^\top(\alpha\ghat) - \frac{L}{2}\|\alpha\ghat\|^2
  = J(\theta_t) + \alpha\,\gtrue^\top\ghat - \frac{\alpha^2 L}{2}\|\ghat\|^2
  \label{eq:before-expectation}
\end{equation}

\textbf{Step 3: Take expectation over $\ghat$.}

\begin{align}
  \E[J(\theta_{t+1})]
  &\geq J(\theta_t) + \alpha\,\E[\gtrue^\top\ghat] - \frac{\alpha^2 L}{2}\E[\|\ghat\|^2]
  \label{eq:exp-step}
\end{align}
We compute the two remaining expectations.

\emph{First expectation:} Since $\E[\ghat] = \gtrue$ (unbiased estimator):
\begin{equation}
  \E[\gtrue^\top\ghat] = \gtrue^\top\E[\ghat] = \gtrue^\top\gtrue = \|\gtrue\|^2
  \label{eq:first-exp}
\end{equation}

\emph{Second expectation:}
We decompose $\ghat = \gtrue + (\ghat - \gtrue)$:
\begin{align}
  \E[\|\ghat\|^2]
  &= \E[\|\gtrue + (\ghat - \gtrue)\|^2] \notag\\
  &= \E[\|\gtrue\|^2 + 2\gtrue^\top(\ghat-\gtrue) + \|\ghat-\gtrue\|^2] \notag\\
  &= \|\gtrue\|^2 + 2\gtrue^\top\underbrace{\E[\ghat-\gtrue]}_{=\,0} + \E[\|\ghat-\gtrue\|^2] \notag\\
  &= \|\gtrue\|^2 + \frac{\sigma_t^2}{\neff}
  \label{eq:second-exp}
\end{align}
where in the last step we used \eqref{eq:mse-neff}.

\textbf{Step 4: Assemble the bound.}

Substituting \eqref{eq:first-exp} and \eqref{eq:second-exp} into \eqref{eq:exp-step}:
\begin{equation}
  \boxed{
  \E[J(\theta_{t+1})] - J(\theta_t)
  \;\geq\;
  \alpha\,\|\gtrue\|^2 - \frac{\alpha^2 L}{2}\!\left(\|\gtrue\|^2 + \frac{\sigma_t^2}{\neff}\right)
  }
  \label{eq:improvement-bound}
\end{equation}
Call the right-hand side $f(\alpha)$.
This is a downward-opening quadratic in $\alpha$:
the $\alpha$ term is positive (driving improvement), and the $\alpha^2$ term is negative (penalizing large steps).

\textbf{Step 5: Maximize over $\alpha$ to find the optimal step size.}

Differentiate $f(\alpha)$ with respect to $\alpha$ and set equal to zero:
\begin{align}
  \frac{d}{d\alpha} f(\alpha)
  &= \|\gtrue\|^2 - \alpha L\!\left(\|\gtrue\|^2 + \frac{\sigma_t^2}{\neff}\right)
  = 0 \notag\\
  \Rightarrow\quad
  \alpha
  &= \frac{\|\gtrue\|^2}{L\!\left(\|\gtrue\|^2 + \sigma_t^2/\neff\right)}
  \label{eq:alphaopt-derivation}
\end{align}
Since $f''(\alpha) = -L(\|\gtrue\|^2 + \sigma_t^2/\neff) < 0$, this critical point is a maximum.
The \textbf{optimal step size} is:
\begin{equation}
  \boxed{
  \alphaopt
  = \frac{\|\gtrue\|^2}{L\!\left(\|\gtrue\|^2 + \sigma_t^2/\neff(t)\right)}
  }
  \label{eq:alphaopt-full}
\end{equation}

\subsection{Monotone dependence on $\neff$: full algebra}
\label{sec:monotone}

\textbf{Step 1: Verify $\alphaopt$ is increasing in $\neff$.}

Treat $\alphaopt$ as a function of $\neff$ (all other quantities held fixed).
Write $A = \|\gtrue\|^2$, $S = \sigma_t^2$, $C = L$.
Then:
\begin{equation}
  \alphaopt = \frac{A}{C(A + S/\neff)}
\end{equation}
Taking the derivative with respect to $\neff$:
\begin{equation}
  \frac{d\alphaopt}{d\neff}
  = \frac{A \cdot C \cdot S/\neff^2}{C^2(A + S/\neff)^2}
  = \frac{A S}{C \neff^2 (A + S/\neff)^2} > 0
  \label{eq:monotone}
\end{equation}
This is strictly positive (since $A, S, C, \neff > 0$), confirming that $\alphaopt$ is strictly increasing in $\neff$.

\textbf{Step 2: Noise-dominated regime.}

When the gradient signal is small relative to the estimation noise,
i.e., $\sigma_t^2/\neff \gg \|\gtrue\|^2$, we can drop the $\|\gtrue\|^2$ term in the denominator:
\begin{align}
  \alphaopt
  &= \frac{\|\gtrue\|^2}{L\!\left(\|\gtrue\|^2 + \sigma_t^2/\neff\right)} \notag\\
  &\approx \frac{\|\gtrue\|^2}{L \cdot \sigma_t^2/\neff} \notag\\
  &= \frac{\|\gtrue\|^2 \cdot \neff}{L \cdot \sigma_t^2}
  \label{eq:alphaopt-noisedom}
\end{align}

\textbf{Step 3: Introduce Assumption~\ref{ass:nsr} for a cleaner expression.}

\begin{assumption}[Slowly-varying noise-to-signal ratio]
  \label{ass:nsr}
  The ratio $\sigma_t^2 / \|\gtrue\|^2$ is approximately constant throughout training:
  \[
    \frac{\sigma_t^2}{\|\gtrue\|^2} \approx C_0
  \]
  for some problem-dependent constant $C_0 > 0$.
\end{assumption}

Under Assumption~\ref{ass:nsr}, substituting $\sigma_t^2 = C_0 \|\gtrue\|^2$ into \eqref{eq:alphaopt-noisedom}:
\begin{align}
  \alphaopt
  &\approx \frac{\|\gtrue\|^2 \cdot \neff}{L \cdot C_0 \|\gtrue\|^2} \notag\\
  &= \frac{\neff}{L \cdot C_0}
  \label{eq:alphaopt-linear}
\end{align}
Since $L$ and $C_0$ are constants (with respect to $t$), this gives:
\begin{equation}
  \boxed{\alphaopt \propto \neff(t)}
  \label{eq:alphaopt-prop}
\end{equation}
The optimal learning rate decreases \emph{linearly} with the effective sample size.

\begin{remark}
  Assumption~\ref{ass:nsr} is the key non-trivial hypothesis.
  It requires that both $\|\gtrue\|^2$ (signal) and $\sigma_t^2$ (per-sample variance) decrease
  at similar rates as the policy concentrates.
  This is plausible in the early-to-mid training regime when the policy has not yet converged.
  Near convergence ($\gtrue \approx 0$), both quantities vanish and the assumption becomes less meaningful---
  but at that point $\alphaopt \to 0$ regardless, so the qualitative conclusion holds.
\end{remark}

\begin{remark}
  Note that the relationship is \emph{linear} ($\alphaopt \propto \neff$), not square-root.
  An earlier (incorrect) derivation based on a signal-to-noise ratio condition
  $\|\gtrue\| \geq \|\ghat - \gtrue\|$ gave $\alphaopt \propto \sqrt{\neff}$;
  however, that condition constrains $\|\gtrue\| / (\sigma/\sqrt{n})$, from which $\alpha$
  cancels entirely---it does not bound $\alpha$.
  The correct approach, using the per-step improvement bound, gives the linear relationship.
\end{remark}

\subsection{The supervised learning baseline}

In supervised learning, the dataset $\mathcal{D}$ is fixed.
The sampling distribution does not change with $\theta$,
so $\neff = n$ and $\sigma_t^2$ remain constant throughout training.
Equation~\eqref{eq:alphaopt-full} then gives a fixed $\alphaopt$, consistent with
the standard practice of using a constant or mildly decaying learning rate.
The systematic degradation of $\neff(t)$ studied here has no analogue in supervised learning.

% =========================================================
\section{The On-Policy Feedback Loop}

\subsection{The structural difference}
\label{sec:structural}

In supervised learning, the data distribution and $\neff$ are fixed external quantities.
In on-policy RL, updating $\theta$ changes $\pith$, which changes the sampling distribution,
which changes $\neff(t)$:
\begin{equation}
  \theta_t
  \;\xrightarrow{\;+\alpha\,\ghat\;}\;
  \theta_{t+1}
  \;\Rightarrow\;
  H(\pith[t+1]) < H(\pith)
  \;\Rightarrow\;
  \neff(t+1) < \neff(t)
  \;\Rightarrow\;
  \alpha^*(t+1) < \alphaopt
  \label{eq:loop-chain}
\end{equation}
The middle implication holds because successful RL training concentrates the policy onto
high-reward trajectories, which reduces entropy.
The rightmost implication follows from \eqref{eq:alphaopt-prop}.

This means $\alphaopt$ is a \emph{time-varying}, \emph{monotonically decreasing} quantity,
while fixed-LR training holds $\alpha$ constant.

\subsection{The mismatch condition}
\label{sec:mismatch}

Let $t^*$ be the first step at which $\alpha > \alpha^*(t^*)$.
Since $\alphaopt \propto \neff(t) \to 0$ as $t \to \infty$ (policy fully concentrates),
such a step must eventually arrive under fixed $\alpha > 0$.

Once $\alpha > \alpha^*(t)$, the bound \eqref{eq:improvement-bound} is no longer positive:
the expected per-step improvement is negative, i.e., the expected next-iterate is \emph{worse}
than the current iterate.

\begin{equation}
  \alpha > \alphaopt
  \;\Rightarrow\;
  \E[J(\theta_{t+1})] < J(\theta_t)
  \label{eq:mismatch}
\end{equation}
An update in this regime is dominated by estimation error: $\alpha(\ghat - \gtrue)$ is larger
in magnitude than $\alpha\gtrue$, so the parameter moves away from the true gradient direction.

\begin{remark}[Self-reinforcing collapse---empirical observation, not part of core argument]
  \label{rem:self-reinforce}
  In practice, the degradation after $t^*$ tends to be self-accelerating.
  A noise-dominated update reinforces whichever trajectories happened to appear most in the sampled batch,
  further concentrating $\pith$ and reducing $\neff(t+1)$, which decreases $\alpha^*(t+1)$
  further, increasing the mismatch:
  \begin{equation}
    \alpha > \alpha^*(t)
    \;\Rightarrow\;
    \text{noise-dominated update}
    \;\Rightarrow\;
    \neff(t+1) \ll \neff(t)
    \;\Rightarrow\;
    \alpha \gg \alpha^*(t+1)
    \;\Rightarrow\;
    \cdots
    \label{eq:feedback-loop}
  \end{equation}
  This is consistent with the sharp ``rise-then-fall'' collapse pattern observed experimentally
  (e.g., \texttt{sync\_lr1e-5} peaks near step 50 then drops sharply).

  However, this self-reinforcing property is \emph{not} a logical step in the core argument.
  The mismatch condition~\eqref{eq:mismatch} alone is sufficient to conclude that updates
  after $t^*$ are expected to be harmful.
  The self-reinforcing loop explains the \emph{severity and speed} of collapse, not its existence.
\end{remark}

% =========================================================
\section{Path Quality and Final Performance}

\subsection{What determines the training path}

The sequence of updates $\theta_0 \to \theta_1 \to \cdots \to \theta_T$ traces a path
through parameter space.
The quality of each step is determined by whether $\alpha \leq \alphaopt$ at that step.
\begin{itemize}
  \item When $\alpha \leq \alphaopt$: expected improvement is positive; the path moves reliably toward higher-reward regions.
  \item When $\alpha > \alphaopt$: expected improvement is negative; the path moves unreliably.
\end{itemize}

\subsection{Why the same terminal entropy implies a path-dependent outcome}

All stable experiments in our study converge to similar final entropy ($H \approx 0.07$--$0.10$ nats),
yet AIME accuracy at step 300 ranges from 0.017 to 0.325---a 20$\times$ spread.

Since the terminal entropy is approximately equal, all policies are similarly concentrated
at the end of training.
The difference must lie in \emph{which} mode the policy converged to---that is, in the path.

\begin{center}
\begin{tabular}{lrrl}
\toprule
Experiment & Entropy@300 & AIME@300 & Pattern \\
\midrule
\texttt{sync\_lr3e-6}  & 0.073 & 0.325 & Monotone increase \\
\texttt{sync\_lr5e-6}  & 0.068 & 0.294 & Slight fluctuation \\
\texttt{sync\_lr1e-6}  & 0.078 & 0.281 & Slow increase \\
\texttt{cosine\_floor} & 0.067 & 0.273 & Stable \\
\texttt{sync\_lr1e-5}  & 0.068 & 0.183 & Rise then fall \\
\texttt{clip0.1}       & 0.099 & 0.017 & Delayed catastrophic \\
\bottomrule
\end{tabular}
\end{center}

The path determines which modes are reached before $\neff$ collapses.
Once $\neff \approx 1$, the policy is effectively locked: unreliable updates cannot explore
other regions because the gradient signal about them has vanished entirely.

\begin{proposition}[Path determines outcome]
  \label{prop:path}
  Given a reward landscape with multiple local optima, the quality of the final policy
  is determined by which local optimum the training path reached before $\neff(t)$ became
  too small for reliable updates ($\alpha > \alphaopt$).
\end{proposition}

% =========================================================
\section{The Research Gap}

\subsection{What existing methods do and do not address}

\begin{center}
\begin{tabular}{lll}
\toprule
\textbf{Method} & \textbf{What it does} & \textbf{Gap} \\
\midrule
Fixed LR & $\alpha$ constant & Ignores $\neff(t) \downarrow$ entirely \\
Cosine decay & $\alpha \downarrow$ on fixed schedule & Schedule is problem-independent \\
PPO clipping & Limits per-step ratio & Does not track $\neff$; blind to unsampled actions \\
Entropy bonus & Resists $H(\pith) \downarrow$ & Opposes training objective; no correct fixed $\beta$ \\
Entropy-adaptive LR & $\alpha \propto H(t)/H(0)$ & Uses entropy as proxy for $\neff$; not yet tuned \\
\bottomrule
\end{tabular}
\end{center}

No existing method \emph{directly tracks} $\neff(t)$ and sets $\alpha$ accordingly.
The derivation in Section~\ref{sec:monotone} establishes that the correct scaling is $\alpha \propto \neff$
(linear), but all open-loop methods use problem-independent schedules that may not match
the actual $\neff(t)$ trajectory.

\subsection{The research question}

\begin{rqbox}
\textbf{Research Question.}
In on-policy RL for LLMs, the effective sample size $\neff(t)$ decreases systematically
as the policy concentrates.
This decrease implies a decreasing optimal learning rate $\alphaopt \propto \neff(t)$ (linear).
Current training algorithms do not adapt their update scale to $\neff(t)$,
causing the learning rate to exceed $\alphaopt$ from some step $t^*$ onward,
leading to expected-negative updates that degrade training path quality
and lead to suboptimal final performance.

\medskip
\textbf{How should training algorithms track and adapt to the evolution of $\neff(t)$
throughout on-policy RL training, in order to maintain reliable gradient updates,
preserve training path quality, and improve final performance?}
\end{rqbox}

\subsection{Solution space}

This research question is intentionally solution-agnostic.
Candidate approaches include, but are not limited to:

\begin{enumerate}[noitemsep]
  \item \textbf{Adaptive learning rate}: set $\alpha(t) \propto \hat{\neff}(t)$
        where $\hat{\neff}$ is estimated from rollout diversity or entropy.
  \item \textbf{Adaptive rollout count}: increase $n$ as $H(\pith)$ decreases,
        maintaining approximately constant $\neff$ despite increasing concentration.
  \item \textbf{Temperature-separated sampling}: generate rollouts at higher temperature
        to maintain diversity while computing advantages under the true $\pith$;
        this partially decouples $\neff$ from $H(\pith)$.
  \item \textbf{Diversity-preserving replay}: retain high-diversity samples from
        earlier training steps and mix them with current rollouts,
        effectively boosting $\neff$ without requiring more environment interactions.
  \item \textbf{Distributional trust region}: directly constrain the distributional
        divergence between consecutive policies on the \emph{full} distribution
        (not just sampled actions), preventing the unsampled-action blind spot
        that causes PPO clipping to underestimate policy change.
\end{enumerate}

Each approach addresses the same root cause---$\neff$ degradation---through a different mechanism.
The framework predicts that any method that successfully maintains $\neff$ at a level
where $\alpha \leq \alphaopt$ will produce better training paths and higher final performance.

% =========================================================
\section{Open Questions}
\label{sec:open}

\begin{enumerate}
  \item \textbf{Formal characterization of $\neff(t)$.}
    Definition~\ref{def:neff} uses $\neff(t) = \E[K_n] \approx \min(n, \exp(H(\pith)))$.
    The perplexity $\exp(H)$ is convenient and differentiable, but it is a response-level quantity.
    In practice we log per-token entropy, which aggregates over template tokens, reasoning tokens,
    and conclusion tokens differently.
    Is per-token entropy a reliable proxy for response-level perplexity?
    Or should response-level diversity metrics (unique response fraction) be used instead?

  \item \textbf{Tightness of the linear relationship.}
    The derivation of $\alphaopt \propto \neff$ (linear) in Section~\ref{sec:monotone}
    relies on two conditions:
    (a) the noise-dominated regime ($\sigma_t^2/\neff \gg \|\gtrue\|^2$) and
    (b) Assumption~\ref{ass:nsr} (constant noise-to-signal ratio).
    A tight bound would require characterizing how $\|\gtrue\|^2$ and $\sigma_t^2$ evolve
    with $H(\pith)$---both are problem-dependent.
    Under what distributional conditions does Assumption~\ref{ass:nsr} hold, and how sensitive
    is the conclusion to violations of this assumption?

  \item \textbf{Separating the ceiling problem.}
    The framework predicts that matching $\alpha$ to $\alphaopt$ improves path quality.
    Whether path quality improvements can push past the 32.5\% AIME ceiling
    depends on whether that ceiling is due to path quality (solvable within this framework)
    or model capacity and reward signal sparsity (fundamental limits).
    New experiments with better-tuned adaptive methods are needed to distinguish these.

  \item \textbf{Optimality of different proxy metrics.}
    Entropy is a coarse proxy for $\neff$.
    Response-level diversity metrics (unique response fraction, pairwise response distance)
    are more direct but harder to compute and differentiate.
    What is the minimal observable that suffices for reliable $\neff$ tracking?
\end{enumerate}

% =========================================================

\begin{thebibliography}{9}

\bibitem{casella2002statistical}
G.~Casella and R.~L. Berger.
\newblock \emph{Statistical Inference}, 2nd ed.
\newblock Duxbury, 2002.
\newblock (Theorem 5.2.4: variance of sample mean for i.i.d.\ vector observations;
  Theorem 4.3.1: covariance decomposition.)

\bibitem{bottou2018optimization}
L.~Bottou, F.~E. Curtis, and J.~Nocedal.
\newblock Optimization methods for large-scale machine learning.
\newblock \emph{SIAM Review}, 60(2):223--311, 2018.
\newblock (Lemma A.1: descent lemma for $L$-smooth objectives;
  Equation (4.8): per-step improvement bound for SGD with noisy gradients.)

\end{thebibliography}

\end{document}
